% Generated draft: requires mathematical and source review.
\documentclass[11pt]{article}
\usepackage[margin=25mm]{geometry}
\usepackage{fontspec}
\setmainfont{Latin Modern Roman}
\usepackage{amsmath,amssymb,mathtools,unicode-math}
\setmathfont{Latin Modern Math}
\usepackage{enumitem,needspace,xurl,hyperref}
\setlength{\parindent}{0pt}
\setlength{\parskip}{5pt}
\setlength{\emergencystretch}{3em}
\newcommand{\N}{\mathbb{N}}
\newcommand{\Z}{\mathbb{Z}}
\newcommand{\Q}{\mathbb{Q}}
\newcommand{\R}{\mathbb{R}}
\newcommand{\C}{\mathbb{C}}
\newcommand{\F}{\mathbb{F}}
\newcommand{\A}{\mathbb{A}}
\newcommand{\PP}{\mathbb{P}}
\newcommand{\E}{\mathbb{E}}
\newcommand{\eps}{\varepsilon}
\newcommand{\dd}{\,\mathrm{d}}
\begin{document}
\section*{UnsolvedMath ID: 800411}
\textbf{Source problem code:} catalogue entry
\subsection{Definitions and mathematical statement}
Let \(\C_v\) denote the valued field used in the source, with absolute value \(|\cdot|_v\), and write \(\mathsf{P}^{1}_{\mathrm{Berk}}\) for its Berkovich projective line. The supplied excerpt does not specify the construction of \(\C_v\), its characteristic, or its valuation hypotheses. These ambient hypotheses must be recovered from the source before the conjecture's full scope can be fixed; the statement below retains the source's notation and intended non-Archimedean setting.

Let \(f\in\C_v(z)\) have degree at least \(2\), and let \(\mathcal J_f\subseteq\mathsf{P}^{1}_{\mathrm{Berk}}\) be its Berkovich Julia set. Type I points are the classical points \(\PP^1(\C_v)\). For a periodic point \(x\in\PP^1(\C_v)\) of least period \(n\geq1\), define its multiplier by
\[
\lambda_f(x)=\left.\frac{\dd}{\dd t}\bigl(t\circ f^{\circ n}\circ t^{-1}\bigr)\right|_{t=0},
\]
where \(t\) is a local coordinate with \(t(x)=0\). This definition also applies at infinity and is independent of the chosen local coordinate. The point is repelling if \(|\lambda_f(x)|_v>1\). Set
\[
R_f^{\mathrm I}=\bigl\{x\in\PP^1(\C_v):x\text{ is periodic under }f\text{ and }|\lambda_f(x)|_v>1\bigr\}.
\]
The supplied Conjecture 34 asserts that
\[
\mathcal J_f\cap\PP^1(\C_v)
=
\overline{R_f^{\mathrm I}}^{\,\PP^1(\C_v)}.
\]
Here the closure is taken within the classical projective line, using its usual topology induced by \(|\cdot|_v\). The excerpt does not explicitly name this topology; this is the standard interpretation used in this draft.
\subsection{Short English statement}
In the source's intended non-Archimedean setting, the classical points of the Berkovich Julia set are conjecturally exactly the limits, within the classical projective line, of repelling classical periodic points. The excerpt leaves the precise ambient field hypotheses unspecified.
\subsection{Sources}
\begin{itemize}
\item Local UnsolvedMath registry, permanent ID \(800411\), exact source problem code \texttt{catalogue entry}; supplied title: \emph{Open Problems in Arithmetic Dynamics --- Conjecture 34}.
\item Benedetto, DeMarco, Ingram, Jones et al., \emph{Current Trends and Open Problems in Arithmetic Dynamics} (2018), Conjecture 34. Supplied source URL: \url{https://arxiv.org/abs/1806.04980}. The registry attributes the proposal to these authors in 2018.
\item Supplied provenance records access on 2026-10-16 and extraction by \texttt{arxiv\_latex}. These are supplied provenance details, not an assertion that this draft independently visited or verified the source.
\item The supplied registry labels the problem \texttt{open}; that label is source data, not an independently verified literature assessment.
\end{itemize}
\end{document}
